%%%%%%%%%%  scan idx 79  |  folio 66  %%%%%%%%%%
La
formule (iv) exprime que $i_!$ est un homomorphisme d'anneaux si l'on introduit sur $\underline K(Y)$ une nouvelle structure multiplicative, par la définition $x_N x'=xx'\lambda_{-1}(N)$ (cf.\ Chapitre I, \S 4). Elle se complète par la formule~:
\begin{enumerate}[label=(\roman*)]
\setcounter{enumi}{4}
\item $i_!(\lambda^p(N,x))=\lambda^p(i_!(x))$ pour $p\geq1$,
\end{enumerate}
qui signifie que $i_!$ définit un $\lambda$-homomorphisme de $\underline K(Y)_N$ dans $\underline K(X)$. Cette formule (v), certainement vraie dans tous les cas, n'est démontrée pour l'instant qu'en caractéristique $0$, car elle résulte de la conjonction des théorèmes 1.4 et 2.8.

On a des formules analogues pour $i_*$ et $i^*$~:
\begin{enumerate}[label=(\roman* bis)]
\item $i^*$ est un homomorphisme d'anneaux gradués.
\item $i_*(x.i^*(y))=i_*(x).y$, d'où en particulier $i_*(i^*(y))=\xi' y$, avec $\xi'=i_*(1)$,
\item $i^*(i_*(x))=x.(-1)^q N^q$, d'où en particulier $i^*(\xi')=(-1)^q N^q$,
\item $i_*(x)i_*(x')=i_*(xx'(-1)^q N^q)$,
\item $i_*$ augmente les degrés de $q$ unités.
\end{enumerate}

Dans ces formules, $q$ est la codimension de $Y$ dans $X$, et l'on pose $N^i=C^i(N)$. Seules les formules (iii bis) et (iv bis) demandent une démonstration~; j'ignore si elles sont bien connues (et même si elles sont vraies~!!!) et me borne à noter qu'elles résultent par réduction des formules (ii) et (iv), \emph{pourvu qu'on se place dans} $G(\underline K(Y))$ \emph{et} $G(\underline K(X))$ \emph{et non dans} $A(Y)$ \emph{et} $A(X)$. On utilise alors le fait que $i_*$ et $i^*$ proviennent par passage au quotient de $i_!$ et $i^!$, en notant que $i_!$ décale la filtration de $q$ unités.

